\justifying
\NSFCBodyText

\subsubsection{项目特色}

本项目不把中文语料质量研究处理成“再训练一个评分器”或“寻找更优过滤阈值”，而是把相近低质量表征背后的生成原因作为研究对象。结构破损与结构完整但模型低熟悉的中文语体，可能在困惑度、规范性或高质量来源相似度上获得近似判断；本项目利用中文语体与书写异质性构造确认性反例，并以可逆修复、特异支持恢复和伪恢复检验两类样本的训练效用是否来自不同路径。结果即使不产生更高的平均榜单分，只要能支持或否定路径假设，仍可形成关于语料质量表征含义的科学结论。

项目还把失败边界纳入研究设计。结构或支持路径未通过负对照时即删除相应机制，独立来源中出现反号或校准失败时明确拒绝外推，不通过增加语体、阈值和模型寻找有利结果。这套分阶段判定规则使“路径是否存在、何时成立、何时不能迁移”具有同等证据地位。

\subsubsection{创新点}

\paragraph{（1）从质量—覆盖相关权衡推进到低质量表征的路径特异分解。}现有工作已认识到强过滤可能损失多样性，本项目不重复这一结论，而是用保持语义支持的结构完整/破损与 sham 对照、同语体非冗余补回、同语体窄支持补回和主流文本伪补回定义两套反事实，分别识别结构缺陷损伤及非冗余语义支持损失。两套估计量不被直接解释为同一种过滤效应的方向比较，其路径特异性由负对照和等效区间判定，不依赖一个可任意加权的综合质量分。

\paragraph{（2）检验中文表面规范性与训练效用的非等价机制。}以语义匹配的规范书面表达和结构完整但模型低熟悉的中文语体为确认性对象，比较至少两类异质质量表征，检验模型低熟悉度是否被误作数据低质量；再以能力定向梯度对齐和独立能力终点区分其与真实结构破损的学习后果。中文在这里是产生“评分相近、因果来源不同”压力测试的理论条件，而非更换数据集的本地化标签。

\paragraph{（3）揭示双路径效应的非对称外部边界。}结构路径重点检验效应方向能否在独立来源保持，支持路径重点检验其随处理前知识状态和目标支持重叠变化的条件效应。来源 A 的协议、效应方向和区间在进入来源 B 前冻结，以“直接运输—条件校准—拒绝外推”描述结果。创新落在两条路径具有何种不同的外部效度结构，校准和拒绝规则仅为检验工具，不被包装成新的通用算法。
